% Supplementary methods and results.
\appendix
\setcounter{figure}{0}
\setcounter{table}{0}
\renewcommand{\thefigure}{A\arabic{figure}}
\renewcommand{\thetable}{A\arabic{table}}
\renewcommand{\theHfigure}{appendix.\arabic{figure}}
\renewcommand{\theHtable}{appendix.\arabic{table}}

\section{Model comparison and vocabulary interpretation}\label{app:models}

Probe AUC measures discrimination between passing and failing programs. Vocabulary decoding identifies output tokens aligned with the probe direction.

The layerwise reconstructions peak at different layers from the original selections in five of six models. Table~\ref{tab:model_peaks} reports the original estimates separately from the curves in Figure~\ref{fig:depth}. Final-normalization conventions also differ; Appendix~\ref{app:final_layer} compares raw and normalized final-block states directly.

\begin{table}[htbp]
\caption{Peak estimates from the original six-model comparison.}
\label{tab:model_peaks}
\centering\small
\begin{tabular}{@{}llcc@{}}
\toprule
Model & Params & Original layer & Original AUC\\
\midrule
LLaDA-flash          & 102B         & L20 & 0.81\\
DiffuCoder-7B        & 7.6B         & L22 & 0.78\\
Dream-Coder-7B       & 7.6B         & L28 & 0.76\\
LLaDA2.0-mini        & 1.4B active  & L14 & 0.74\\
UltraLLaDA-8B        & 8B           & L27 & 0.72\\
Stable-DiffCoder-8B  & 8.25B        & L19 & 0.67\\
\bottomrule
\end{tabular}
\end{table}

The LLaDA2-MoE models show stronger enrichment of wrongness terms than the dense coders despite comparable correctness readouts (Figure~\ref{fig:interpretable}).

\begin{figure}[htbp]
\centering
\includegraphics[width=.64\linewidth]{jspace_readable_vs_interpretable.pdf}
\caption{Probe AUC (horizontal axis) and wrongness-word enrichment (vertical axis). All six models have above-chance readouts, but only the LLaDA2-MoE pair strongly enriches explicit ``wrong / incorrect'' words. The dense coders' vocabulary enrichment is near the random-direction baseline.}
\label{fig:interpretable}
\end{figure}

\clearpage
\section{Capture protocol, scoring, and uncertainty}\label{app:methods}

Table~\ref{tab:cohorts} lists the analysis samples. Within-problem AUC uses tasks containing both passing and failing attempts. The family comparison and follow-up experiments use separate captures.

\begin{table}[htbp]
\centering\small
\caption{Problems and examples included in each analysis.}
\label{tab:cohorts}
\begin{tabular}{@{}p{.49\linewidth}rr@{}}
\toprule
Cohort & Problems & Examples\\
\midrule
DiffuCoder, 32 denoising steps & 19 & 95\\
DiffuCoder, 16 denoising steps & 54 & 324\\
Qwen2.5-Coder, generated attempts & 56 & 448\\
DiffuCoder, MBPP+ mutants & 52 & 392\\
DiffuCoder, HumanEval+ mutants & 44 & 334\\
LLaDA-flash / Stable remeasurement & 53 & 399 each\\
\bottomrule
\end{tabular}
\end{table}

Before mixed-problem filtering, the 32-step DiffuCoder sample contains 120 tasks with five attempts each, and the Qwen2.5-Coder sample contains 180 with eight attempts each. Analysis uses DiffuCoder's last-step response-token mean and Qwen's completion-token mean.

Activations are centered within problem. Five folds are assigned by problem using NumPy seed 0; attempts from the same problem never cross folds. The normalized difference between passing and failing training activations defines each fold's direction. AUCs of held-out projections are averaged equally across mixed problems. The implementation uses sorted ranks without explicit tie averaging. The maximum of the out-of-fold curve selects the reported layer, without a separate layer-selection holdout.

Diffusion confidence controls use mean visible-token log probability, negative mean visible-token entropy, and mean completion log probability with response positions masked. Diffusion logits at $i$ score token $i$; autoregressive logits score the following token. Visible and masked passes condition on different information.

\paragraph{Null distributions.}\label{app:uncertainty} We use 200 within-problem label permutations. For the confidence comparison, each permutation evaluates fixed out-of-fold scores and repeats layer selection; probe directions are not refitted. The perturbation analyses evaluate fixed scores at a preselected layer. Final-layer analyses select the peak before computing its null, so those peak scores do not account for layer selection.

\paragraph{Intervals.} The probe--confidence bootstrap implementation merges repeated problem identifiers after resampling, losing their sampling multiplicity. We therefore do not interpret its intervals as cluster-bootstrap confidence intervals. Correct intervals require resampling paired per-problem AUC differences, with layer selection addressed separately.

\clearpage
\section{Perturbation controls}\label{app:perturbations}

DiffuCoder, LLaDA-mini, and LLaDA-flash use noise seeds $s=0,1,2,3,4$, mapped to NumPy seeds $1000000+s$. One random unit direction per example is shared across positions and scaled by the mean token-state norm. The remaining three models were evaluated with one noise realization. Figure~\ref{fig:noise} shows the three models with repeated seeds.

\begin{figure}[htbp]
\centering
\includegraphics[width=.60\linewidth]{g1_param_correctness_noise.pdf}
\caption{Mean correctness-probe AUC across five noise realizations. Noise of relative magnitude $m$ is added at the probed layer. Performance declines at different rates across the three models.}
\label{fig:noise}
\end{figure}

\begin{figure}[htbp]
\centering
\includegraphics[width=.62\linewidth]{auc_vs_causal.pdf}
\caption{Correctness-probe AUC after attention-edge removal. Most sub-chance dips and subsequent increases have per-point $|z|<2$ relative to their shuffled-label nulls. The curves use a single seed.}
\label{fig:causal}
\end{figure}

\clearpage
\section{Steering configuration and results}\label{app:steering}

The follow-up steering experiments intervene at DiffuCoder block index 18. For a unit direction $v$, we add $\alpha\bar r v$ to selected response positions at every denoising step, where $\bar r$ is the row's mean response-token residual norm and $\alpha=\pm0.4$. Global arms affect all response positions; conditional arms affect only positions still containing the mask token. All eleven arms share one generation batch: an untreated baseline, probe directions, random directions, and a learned direction.

A fresh random unit direction is sampled per task with NumPy seed $10000+i$, where $i$ is the evaluation-local task index. The evaluation configuration specifies 256 new tokens, 32 denoising steps, and temperature 0. Repeated untreated rows in the earlier bf16 sweep had different pass counts.

The learned direction starts from the probe and minimizes reference-token negative log probability from a masked-response input while model weights are frozen. Adam runs for 25 steps at learning rate 0.005, alternating between two groups of four references.

Table~\ref{tab:steerarms} reports results on 30 tasks. The learned direction has cosine similarity 0.108 with the probe direction. Its positive global intervention reduces passing programs from 15 to 5 and compiling programs from 26 to 9.

\begin{table}[htbp]
\centering\small
\caption{Passing and compiling programs out of 30 evaluation tasks.}
\label{tab:steerarms}
\begin{tabular}{@{}lrr@{}}
\toprule
Intervention & Pass tests & Compile\\
\midrule
Baseline & 15 & 26\\
Probe, global $+$ / $-$ & 14 / 5 & 24 / 16\\
Probe, conditional $+$ / $-$ & 15 / 10 & 25 / 21\\
Random, global $+$ / $-$ & 15 / 14 & 22 / 24\\
Random, conditional $+$ & 16 & 25\\
Learned, global $+$ / $-$ & 5 / 15 & 9 / 27\\
Learned, conditional $+$ & 7 & 12\\
\bottomrule
\end{tabular}
\end{table}

The main steering figure summarizes a separate sweep; its error bar is the standard deviation across ten random directions.

\clearpage
\section{Final-layer measurements and the compression hypothesis}\label{app:final}\label{app:final_layer}

A transformer block's raw output and its final normalized state are different measurements. Their inconsistent treatment in the initial captures motivated separate measurements before and after final normalization. We also tested whether projection toward output-token directions could explain a late decline in probe AUC (Figure~\ref{fig:compression}).

\begin{figure}[htbp]
\centering
\resizebox{0.76\textwidth}{!}{%
\begin{tikzpicture}[font=\small]
  \draw[-{Stealth[length=2mm]}, soft, line width=0.9pt] (0,0) -- (11.4,0);
  \draw[-{Stealth[length=2mm]}, soft, line width=0.9pt] (0,0) -- (0,3.4);
  \node[soft, below] at (5.3,-0.1) {network depth (layer) $\rightarrow$};
  \node[soft, rotate=90, anchor=south, align=center] at (-0.45,1.7) {readability of the\\correctness signal};
  \draw[soft, dashed, line width=0.7pt] (0,0.45) -- (10.4,0.45);
  \node[soft, font=\scriptsize, above] at (1.1,0.45) {chance};
  \fill[emph!9] (8.7,0) rectangle (10.4,3.2);
  \node[emph, font=\scriptsize, align=center] at (9.55,-0.42) {final\\layer(s)};
  \draw[accent, line width=1.5pt] plot[smooth] coordinates
    {(0.2,0.47)(1,0.95)(2,1.3)(3.2,1.75)(4.5,2.2)(6,2.5)(7.2,2.55)(8,2.4)(8.8,1.75)(9.5,1.0)(10.15,0.62)};
  \fill[accent] (7.2,2.55) circle (2.2pt);
  \node[accent, font=\footnotesize, align=center] at (3.7,2.7) {built with depth\\(readable)};
  \node[emph, font=\footnotesize, align=center] at (9.55,2.75) {collapses?};
  \node[draw=ink!45, fill=ink!5, rounded corners=2pt, align=center, inner sep=5pt,
        text width=2.3cm, font=\footnotesize] (tok) at (12.7,1.5) {output token distribution};
  \draw[-{Stealth[length=2.4mm]}, soft, line width=1.0pt] (10.2,0.9) .. controls (11.2,1.2) .. (tok.west);
\end{tikzpicture}}
\caption{Schematic of the proposed compression mechanism, not measured data. A late conversion toward output-token directions might reduce a linear probe's access to correctness information. The DiffuCoder-7B tests below did not support this explanation.}
\label{fig:compression}
\end{figure}

The DiffuCoder final-layer battery uses five-fold splits by problem. Logistic regression uses $C=0.05$ and at most 2000 iterations. The multilayer perceptron has one 64-unit hidden layer, regularization parameter 0.01, at most 600 iterations, and random seed 0. The vocabulary-subspace test uses the leading 64 right singular vectors of the mean-centered output embedding matrix. The normalization test applies RMSNorm to an already pooled final-block vector; this is not generally equivalent to normalizing each token before pooling.

In that capture, the final-block linear AUC is 0.838 against a peak of 0.842; the normalized pooled-vector AUC is 0.816. Splitting the representation by the selected vocabulary subspace gives higher AUC outside it than inside it: 0.829 versus 0.781 at the peak, and 0.811 versus 0.741 at the final block. The tested nonlinear decoder does not yield a higher point estimate than the peak linear probe. DiffuCoder therefore retains a strong correctness read at its final block, without evidence for the proposed compression mechanism.

The LLaDA-flash and Stable remeasurements instead capture final-block token states directly, apply final normalization to each token, and then average. They use compilable reference-code mutants, with passing and failing mutants for each retained base problem. Inputs are truncated to 640 tokens and processed singly. LLaDA-flash has peak AUC 0.8002, raw-final AUC 0.6840, and normalized-final AUC 0.6976. For Stable-DiffCoder, the corresponding AUCs are 0.7301, 0.6086, and 0.6325.

Stable-DiffCoder also shows a final-layer decline. The effect is therefore not confined to the 102B model, and these comparisons do not establish a scale threshold.

\clearpage
\section{Diffusion versus autoregressive models}\label{app:comparison}\label{app:architecture}

Is a readable correctness signal special to diffusion language models? We applied the same type of linear probe to two autoregressive code models on SWE-bench-derived rollouts. One of them, Qwen3-Coder, reads nearly as well as the diffusion models in this comparison, with an AUC of 0.70. The other, GLM-4.5-Air, is much closer to its shuffled-label baseline, with an AUC of 0.56 (Figure~\ref{fig:comparison}).

We later added Qwen2.5-Coder on MBPP+, the programming benchmark used in the diffusion experiments. Its probe reaches an AUC of 0.677, while its own log-probability confidence reaches 0.703. The probe peaks near the end of the layer stack, whereas DiffuCoder's peak appears earlier. Their training histories, task samples, and capture settings differ, so the comparison does not isolate an effect of architecture.

\begin{figure}[htbp]
\centering
\includegraphics[width=.64\linewidth]{compare_ar.pdf}
\caption{Correctness probes across diffusion and autoregressive models. Each model is scored against its own shuffled-label null; the displayed $\sigma$ values describe distances from those separate nulls and should not be compared across bars as a common effect size. Task sources and training also differ.}
\label{fig:comparison}
\end{figure}

During diffusion denoising, a correctness signal can be present while multiple positions in the program remain revisable. In an autoregressive model, the tokens already written are committed, and a read taken mid-generation can only shape what comes next. Diffusion therefore offers an opportunity to use information about the program while the answer can still change. Whether that opportunity is useful depends on when the signal appears and whether an intervention supplies a repair rather than just a correctness verdict.

In the matched confidence comparisons, the probe offers no consistent advantage over the model's own confidence. Nor did positive probe or learned-direction steering improve pass counts in the 30-task evaluation. The tested models provide a readable signal; the interventions leave open how to turn that signal into a useful repair.
